% Einfügen von Tabellen, Abbildungen und Listings
\newcommand{\tabelle}[1]{\input{tables/#1}}
\newcommand{\abbildung}[1]{\input{figures/#1}}
\newcommand{\code}[1]{\input{listings/#1}}

% Gebräuchliche Abkürzungen
\newcommand{\zB}{z.\,B.\@\xspace}      % zum Beispiel
\newcommand{\bzw}{bzw.\@\xspace}       % beziehungsweise
\newcommand{\usw}{usw.\@\xspace}       % und so weiter
\newcommand{\ua}{u.\,a.\@\xspace}      % unter anderem
\newcommand{\uU}{u.\,U.\@\xspace}      % unter Umständen
\newcommand{\dH}{d.\,h.\@\xspace}      % das heißt
\newcommand{\vgl}{vgl.\@\xspace}       % vergleiche
\newcommand{\ca}{ca.\@\xspace}         % circa
\newcommand{\evtl}{evtl.\@\xspace}     % eventuell
\newcommand{\bspw}{bspw.\@\xspace}     % beispielsweise
\newcommand{\sog}{sog.\@\xspace}       % sogenannt
\newcommand{\va}{v.\,a.\@\xspace}        % "vor allem"
\newcommand{\zT}{z.\,T.\@\xspace}        % "zum Teil"
\newcommand{\seite}{S.\@\xspace}    % Seite

% Referenzierung
\newcommand{\anhang}[2]{%
  \cref{#2}%
  ~(Anhang~\ref{#1} auf Seite~\pageref{#2})\@\xspace%
}
